\documentclass[12pt,fleqn]{article}

\usepackage[T2A]{fontenc}
\usepackage{amssymb,amsmath,mathrsfs,amsthm}
\usepackage[russian]{babel}
\usepackage{graphicx}
\usepackage{indentfirst}
\usepackage{lipsum}
% \documentclass{article}
\usepackage{array}
\usepackage{caption}
\usepackage{longtable}
\usepackage[inkscapeformat=png]{svg}
\usepackage{svg}
\usepackage{caption}
\renewcommand{\baselinestretch}{1.5}


\textheight=24cm
\textwidth=16cm
\oddsidemargin=5mm
\evensidemargin=-5mm
\marginparwidth=36pt
\topmargin=-2cm
\footnotesep=3ex
\raggedbottom
\tolerance 3000
\clubpenalty=10000
\widowpenalty=10000

\usepackage[T2A]{fontenc}
\usepackage{amssymb,amsmath,mathrsfs,amsthm}
\usepackage{float}
\usepackage[russian]{babel}
\usepackage{graphicx}
\usepackage{indentfirst}
\usepackage{mdframed}
\usepackage{hyperref}
\usepackage{changepage}
\usepackage{caption}
\captionsetup{font=footnotesize}
\usepackage{multirow}
\usepackage{setspace}
\usepackage{booktabs}
\textheight=24cm
\textwidth=16cm
\oddsidemargin=5mm
\evensidemargin=-5mm
\marginparwidth=36pt
\topmargin=-2cm
\footnotesep=3ex
\raggedbottom
\tolerance 3000
\clubpenalty=10000
\widowpenalty=10000



\begin{document}

\begin{titlepage}
\begin{center}
    Московский государственный университет имени М. В. Ломоносова

    \bigskip
    \includegraphics[width=50mm]{msu.jpeg}

    \bigskip
    Факультет Вычислительной Математики и Кибернетики\\
    Кафедра Математических Методов Прогнозирования\\[10mm]

    Загатин Даниил Ильич\\[10mm]

    {\large\bfseries
        <<Применение графовых нейронных сетей в рекомендательных системах>>
    }\\[10mm]

    
        ВЫПУСКНАЯ КВАЛИФИКАЦИОННАЯ  РАБОТА
    \\[30mm]

    \begin{flushright}
        \parbox{0.5\textwidth}{\raggedleft
            \textbf{Научный руководитель:}\\
            к.ф-м.н., доцент\\
            В.В.Китов
        }
    \end{flushright}

    

    \vspace{\fill}
    Москва, 2026
\end{center}
\end{titlepage}

\newpage
\begin{abstract}

Традиционные подходы к рекомендательным системам часто опираются на неявную обратную связь либо бинаризуют пользовательские оценки, из‑за чего теряется информация о силе предпочтения. В отличие от них, данная работа фокусируется на использовании явной обратной связи (численных оценок) в графовых моделях рекомендаций фильмов. Ключевая идея - задействовать оценку пользователя как дополнительный структурный сигнал в графе: в качестве типа ребра, веса связи или временно‑взвешенного вклада. Дополнительно учитывается признаковое описание фильмов. В рамках исследования базовые методы (матричная факторизация и двубашенная модель) служат точкой отсчёта для сравнения эффективности графовых нейронных сетей и гибридных моделей, интегрирующих графовые представления с результатами двубашенной архитектуры.


\end{abstract}

\newpage
\renewcommand{\contentsname}{Содержание}
\tableofcontents


\newpage
\section{Введение}
\label{sec:introduction}

Рекомендательные системы используются в онлайн-кинотеатрах, музыкальных
платформах, интернет-магазинах и других цифровых сервисах, где пользователю
нужно выбрать несколько объектов из большого каталога. В таких приложениях
важно правильно упорядочить
кандидатов и поместить наиболее подходящие фильмы в верхнюю часть списка.
Поэтому в работе рассматривается задача ранжирования первых $K$ объектов рекомендаций.
Рекомендательные модели решают эту задачу разными способами.
Матричная факторизация строит латентные представления пользователей и фильмов
по разреженной матрице оценок~\cite{koren2009matrix}. Попарные функции потерь,
например BPR~\cite{rendle2009bpr}, напрямую оптимизируют относительный порядок
положительных и отрицательных объектов. Двубашенные модели и последовательные
архитектуры позволяют учитывать историю пользователя и признаки фильмов
~\cite{covington2016youtube,kang2018sasrec}. Однако во многих практических
постановках взаимодействие сводится к бинарному событию: пользователь либо
взаимодействовал с фильмом, либо нет. Если в данных доступны явные рейтинги,
такая бинаризация отбрасывает информацию о силе предпочтения.

В рекомендации фильмов численная оценка несёт важный смысл. Оценки $5$, $3$ и
$1$ не просто фиксируют факт взаимодействия, а отражают разные уровни отношения
пользователя к фильму. При этом прямое предсказание рейтинга не полностью
соответствует прикладной цели, потому что системе нужно построить персонализованный список рекомендаций а не рейтинг лучших фильмов. Отсюда возникает основная проблема: необходимо
использовать значение рейтинга как источник информации, но обучать и оценивать
модель в терминах качества рекомендаций.

Графовые нейронные сети естественно подходят для такой постановки. Историю
оценок можно представить как двудольный граф пользователь-фильм, где вершины
соответствуют пользователям и фильмам, а рёбра соответствуют наблюдаемым
оценкам. В отличие от других, графовые методы распространяют информацию по структуре
соседства и используют связи более высокого порядка. При этом рейтинг можно
учитывать несколькими способами: задавать отдельный тип ребра для каждого
значения оценки, превращать оценку в вес связи, обучать функцию рейтингового
веса или ослаблять старые взаимодействия с помощью временного затухания. Эти
идеи связаны с моделями GCMC, NGCF и LightGCN
~\cite{vandenberg2017gcmc,wang2019ngcf,he2020lightgcn}.

Идея же предлагаемого подхода состоит в том, чтобы сравнить различные способы учёта рейтинга в графовых нейронных сетях и, совместив лучший в ансамбле с сильной двубашенной моделью получить прирост качества за счёт одновременного использования информации из графовой структуры данных и последовательной истории пользователя. Также используются дополнительно сгенерированное признаковое описание объектов, которое нацелено на улучшение инициализации их векторных представлений.

Дальнейший текст работы организован следующим образом:
\begin{itemize}
    \item в разделе~\ref{sec:problem} формализуется задача рекомендации фильмов
    по истории явных оценок. Вводятся множества пользователей и фильмов,
    описывается постановка Топ-$K$ ранжирования, задаётся способ перехода от
    численных рейтингов к релевантности и приводятся метрики качества,
    используемые для сравнения моделей;

    \item в разделе~\ref{sec:literature_methods} рассматриваются методы, на
    которых основано исследование. Сначала описываются базовые подходы к
    рекомендациям: матричная факторизация, попарное обучение ранжированию и
    двубашенные модели. Затем рассматриваются графовые методы для
    рекомендательных систем и способы учёта явных рейтингов в графе: рейтинг
    как тип ребра, рейтинг как вес связи, обучаемая функция рейтингового веса и
    временно-взвешенное распространение сигнала. Формулируется идея ансамблирования различных методов. В таблице ~\ref{tab:models_overview} приведены все рассматриваемые методы.

    \item в разделе~\ref{sec:experiments} описывается экспериментальная часть
    работы. Приводится схема подготовки датасета, временное разбиение
    на обучающую, валидационную и тестовую части, а также сценарии тёплого
    старта и холодного старта по пользователю и фильму. Далее описываются
    реализованные модели, настройки обучения, протокол оценки и результаты
    сравнения матричной факторизации, двубашенной модели, графовых моделей и их
    ансамблей;

    \item в разделе~\ref{sec:conclusion} формулируются основные выводы
    исследования. Отдельно обсуждаются ограничения рассмотренных подходов,
    влияние способов учёта рейтинга на качество рекомендаций и возможные
    направления дальнейшего развития работы.
    
\end{itemize}

\section{Постановка задачи}
\label{sec:problem}

\subsection{Основные обозначения}

Пусть $\mathcal{U}$ обозначает множество пользователей, а $\mathcal{I}$ обозначает множество фильмов. Размеры этих множеств обозначим через $|\mathcal{U}|=n$ и $|\mathcal{I}|=m$. Наблюдаемая история взаимодействий задаётся множеством
$$
    \Omega =
    \{(u,i,r_{ui},t_{ui}) \mid u \in \mathcal{U},\ i \in \mathcal{I}\},
$$
где $r_{ui}$ обозначает явную оценку, поставленную пользователем $u$ фильму $i$, а $t_{ui}$ обозначает момент выставления оценки.

Рейтинг может принимать значения, например, из множества полуцелых оценок
$
    \mathcal{R}=\{0.5,1.0,1.5,\ldots,5.0\}.
$
Далее все формулы записываются для общего конечного упорядоченного множества рейтингов $\mathcal{R}$.

Каждому пользователю $u$ может соответствовать обучаемый вектор
$
    \mathbf{p}_u \in \mathbb{R}^{d},
$
а каждому фильму $i$ соответствует обучаемый вектор
$
    \mathbf{q}_i \in \mathbb{R}^{d}.
$
Здесь $d$ обозначает размерность латентного пространства. Дополнительно фильм может иметь признаковое описание
$
    \mathbf{x}_i \in \mathbb{R}^{f},
$
где $f$ обозначает размерность вектора признаков фильма. В качестве признаков могут выступать жанры, наличие популярных актёров или режиссёров и иные метаданные. Признаки, которые напрямую кодируют будущие пользовательские оценки или содержат информацию из тестовой части, должны быть исключены, чтобы избежать утечки данных.

Историю пользователя до момента времени $t$ обозначим через
$$
    \mathcal{H}_u(t)=
    \{(i,r_{ui},t_{ui}) \in \Omega \mid t_{ui}<t\}.
$$
Модель рекомендаций задаёт функцию оценки релевантности
$$
    s_{\theta}(u,i,t) \in \mathbb{R},
$$
где $\theta$ обозначает параметры модели. Чем больше значение $s_{\theta}(u,i,t)$, тем выше модель ранжирует фильм $i$ для пользователя $u$ в момент времени $t$.

\subsection{Формальная задача ранжирования первых $K$ объектов}

Для каждого пользователя $u$ задано множество кандидатов $\mathcal{C}(u) \subseteq \mathcal{I}$. Из него исключаются фильмы, которые пользователь уже оценивал в обучающей части:
$$
    \mathcal{C}(u)=
    \mathcal{I}\setminus
    \{i \mid (u,i,r_{ui},t_{ui})\in\Omega_{\text{обуч}}\}.
$$
Задача модели состоит в том, чтобы выбрать $K$ фильмов с наибольшими значениями функции оценки релевантности:
$$
    \operatorname{TopK}_{\theta}(u)=
    \underset{i\in\mathcal{C}(u)}{\operatorname{topK}}\,
    s_{\theta}(u,i,t).
$$
Итоговая рекомендация является упорядоченным списком
$$
    R_K(u)=(i_1,i_2,\ldots,i_K),
$$
где
$$
    s_{\theta}(u,i_1,t)\geq
    s_{\theta}(u,i_2,t)\geq
    \ldots\geq
    s_{\theta}(u,i_K,t).
$$

В отличие от задачи предсказания рейтинга, где требуется как можно точнее предсказать численное значение оценки, в задаче ранжирования первых $K$ объектов важен относительный порядок фильмов. Модель может неточно восстанавливать абсолютное значение рейтинга, но при этом хорошо решать задачу рекомендаций, если фильмы с высокой релевантностью располагаются выше нерелевантных фильмов.

\subsection{Бинарная релевантность}

Для оценки качества рекомендаций необходимо определить, какие фильмы считаются релевантными. В бинарной постановке фильм считается релевантным, если его рейтинг не меньше заданного порога $r_{+}$:
$$
    y^{\mathrm{bin}}_{ui}=\mathbb{I}[r_{ui}\geq r_{+}].
$$
Далее в экспериментах используется порог
$
    r_{+}=4.0.
$
Тогда оценки $4.0$, $4.5$ и $5.0$ считаются положительным пользовательским откликом.

Бинарная релевантность удобна для интерпретации и хорошо согласуется с задачей ранжирования первых $K$ объектов, поскольку именно порядок фильмов в выдаче является целевым критерием качества модели. При этом исходный рейтинг не игнорируется, а используется как дополнительный источник информации о силе предпочтения пользователя. Таким образом, он выступает не в роли таргета, а дополнительной основы для построения списка рекомендаций.


\subsection{Проблема сценариев холодного и тёплого старта в системах рекомендаций}

Для более полного анализа качества рассматриваются несколько сценариев оценки,
различающихся тем, насколько хорошо пользователь и фильм представлены в
обучающей выборке.

\textbf{Тёплый старт.} Пользователь и фильм уже встречались в обучающей выборке.
Это основной сценарий коллаборативной фильтрации: модель может использовать как
историю оценок пользователя, так и накопленную информацию о фильме, полученную
из взаимодействий с другими пользователями.

\textbf{Холодный старт по пользователю.} Пользователь отсутствует в основной
обучающей выборке, однако для него доступно небольшое множество начальных
взаимодействий. В этом случае модель должна построить представление пользователя
по ограниченной опорной истории:
$$
    \mathcal{H}^{\mathrm{sup}}_u=
    \{(i,r_{ui},t_{ui})\}_{i=1}^{s_u}.
$$
Такой сценарий проверяет способность модели адаптироваться к новому
пользователю без полного переобучения.

\textbf{Холодный старт по фильму.} Фильм отсутствует в обучающей выборке
взаимодействий. Для него могут быть доступны признаки $\mathbf{x}_i$, но
отсутствует или почти отсутствует коллаборативный сигнал. Этот сценарий
особенно важен для проверки пользы дополнительных признаков фильмов, поскольку
именно они позволяют получить представление нового объекта без достаточной
истории оценок.

\textbf{Полностью холодный старт.} И пользователь, и фильм являются новыми
относительно обучающей выборки. Это наиболее сложный сценарий, в котором модель
должна одновременно использовать начальную историю пользователя и признаковое
описание фильма.

Таким образом, задача работы состоит не только в максимизации качества в
сценарии тёплого старта, но и в проверке устойчивости моделей в условиях
холодного старта, когда информация о взаимодействиях пользователя, фильма или
обоих объектов ограничена.

Следует отметить, что в данной работе количественная оценка проводится для тёплого старта и
холодного старта по пользователю. Сценарии появления новых фильмов не включены
в основной экспериментальный протокол из-за ограничений по объёму данных и
вычислительным ресурсам. Тем не менее они остаются важной практической
постановкой и учитываются при обсуждении ограничений работы и возможных
направлений дальнейшего развития.

\subsection{Метрики качества}
\label{subsec:metrics}

Качество моделей оценивается в задаче построения списка из $K$ рекомендаций. Для каждого пользователя $u$ модель строит упорядоченный список рекомендаций
$$
    R_K(u)=(i_1,i_2,\ldots,i_K).
$$
В экспериментах используются значения $K\in\{20,50\}$.
Множество релевантных фильмов пользователя в тестовой выборке обозначим через $G(u)$. В основной постановке
$$
    G(u)=
    \{i\in\mathcal{I}_{\text{тест}}(u)\mid r_{ui}\geq r_{+}\}.
$$

Основными метриками в работе являются HitRate@20, HitRate@50, nDCG@20 и nDCG@50. Дополнительно в полных таблицах результатов приводятся Recall@20 и Recall@50. Такой выбор согласуется с общепринятой практикой оценки рекомендательных систем и задач ранжирования~\cite{herlocker2004evaluating,gunawardana2009survey,steck2013evaluation}: он позволяет одновременно оценить наличие релевантного фильма в выдаче и качество его позиции внутри ранжированного списка.

\paragraph{HitRate@$K$.}

Метрика HitRate@$K$ равна единице, если в первых $K$ рекомендациях есть хотя бы
один релевантный фильм:
$$
    \operatorname{HitRate}@K(u)=
    \mathbb{I}\left[R_K(u)\cap G(u)\neq\varnothing\right].
$$
Итоговое значение усредняется по пользователям:
$$
    \operatorname{HitRate}@K=
    \frac{1}{|\mathcal{U}_{\text{оцен}}|}
    \sum_{u\in\mathcal{U}_{\text{оцен}}}
    \operatorname{HitRate}@K(u).
$$
Метрика показывает долю пользователей, для которых модель нашла хотя бы один
релевантный объект в первые $K$ позиций списка. В работе используются значения $K=20$ и $K=50$.

\paragraph{Recall@$K$.}

Метрика Recall@$K$ показывает, какая доля релевантных тестовых фильмов пользователя попала в список из $K$ рекомендаций:
$$
    \operatorname{Recall}@K(u)=
    \frac{|R_K(u)\cap G(u)|}{|G(u)|}.
$$
Итоговое значение усредняется по пользователям:
$$
    \operatorname{Recall}@K=
    \frac{1}{|\mathcal{U}_{\text{оцен}}|}
    \sum_{u\in\mathcal{U}_{\text{оцен}}}
    \operatorname{Recall}@K(u).
$$
Recall@$K$ полезен для анализа полноты рекомендаций: если у пользователя есть несколько релевантных фильмов в тестовой выборке, метрика показывает, какую их часть удалось восстановить. Однако Recall@$K$ не учитывает порядок фильмов внутри первых $K$ позиций. Например, если релевантный фильм находится на первой или на двадцатой позиции, вклад в Recall@20 будет одинаковым.

\paragraph{nDCG@$K$.}

Метрика nDCG@$K$ основана на мере дисконтированной накопленной полезности (DCG)~\cite{jarvelin2002dcg} и учитывает не только попадание релевантных фильмов в первые $K$ позиций,
но и их позиции в списке. Сначала вычисляется накопленная полезность с дисконтированием:
$$
    \operatorname{DCG}@K(u)=
    \sum_{j=1}^{K}
    \frac{\operatorname{rel}_{u i_j}}{\log_2(j+1)},
$$
где $\operatorname{rel}_{u i_j}$ обозначает релевантность фильма $i_j$ для
пользователя $u$. В работе используется бинарная релевантность:
$$
    \operatorname{rel}_{ui}=\mathbb{I}[r_{ui}\geq r_{+}].
$$
Нормированная версия метрики имеет вид
$$
    \operatorname{nDCG}@K(u)=
    \frac{\operatorname{DCG}@K(u)}{\operatorname{IDCG}@K(u)},
$$
где $\operatorname{IDCG}@K(u)$ обозначает значение DCG для идеального ранжирования
релевантных фильмов пользователя. Если у пользователя нет релевантных фильмов
в тестовой выборке, он исключается из расчёта пользовательского среднего.

nDCG@$K$ используется как основная метрика качества ранжирования, поскольку
релевантные фильмы должны не только попадать в выдачу, но и располагаться как
можно выше. Для выбранной постановки работы это особенно важно: методы могут находить
одинаковое число релевантных фильмов, но по-разному упорядочивать их внутри
списка из $K$ объектов.

\section{Обзор литературы и рассматриваемые методы}
\label{sec:literature_methods}

В данном разделе рассматриваются основные классы методов, используемых в работе. Для каждого класса сначала
приводится общая идея, затем описывается, каким образом соответствующий подход
используется или модифицируется в данной работе. Основное внимание уделяется
способам учёта явных пользовательских рейтингов в задаче построения списка из $K$ рекомендаций.

\subsection{Матричная факторизация}

Одним из базовых подходов к рекомендательным системам является коллаборативная
фильтрация. Её основная идея состоит в том, что предпочтения пользователя можно
восстановить по поведению похожих пользователей и по структуре уже известных
взаимодействий. Если два пользователя оценивали похожие фильмы, то фильмы,
понравившиеся одному из них, могут быть рекомендованы другому. Такой подход не
требует подробного описания объектов, но существенно зависит от плотности
матрицы взаимодействий.

Классическим примером коллаборативной фильтрации является матричная факторизация
~\cite{koren2009matrix}. Её варианты могут использовать как явную, так и неявную обратную связь~\cite{chen2018matrix}. Пусть
$R\in\mathbb{R}^{n\times m}$ обозначает разреженную матрицу рейтингов, где
$R_{ui}=r_{ui}$, если пользователь $u$ поставил фильму $i$ оценку $r_{ui}$,
и значение отсутствует иначе. Матричная факторизация приближает эту матрицу
произведением двух матриц низкого ранга:
$$
    R\approx P Q^{\top},
$$
где
$$
    P\in\mathbb{R}^{n\times d},
    \qquad
    Q\in\mathbb{R}^{m\times d}.
$$
Строка $\mathbf{p}_u$ матрицы $P$ является представлением пользователя, а строка
$\mathbf{q}_i$ матрицы $Q$ является представлением фильма. Оценка релевантности пары
пользователь--фильм задаётся скалярным произведением:
$$
    s_{\mathrm{MF}}(u,i)=
    b_u+b_i+\langle \mathbf{p}_u,\mathbf{q}_i\rangle,
$$
где $b_u$ и $b_i$ обозначают смещения пользователя и фильма.

Для задачи построения списка из $K$ рекомендаций матричную факторизацию удобно обучать не на
среднеквадратичную ошибку предсказания рейтинга, а на попарную функцию потерь
ранжирования. В частности, используется BPR-подход~\cite{rendle2009bpr}, в
котором модель обучается располагать положительные объекты выше негативных:
$$
    \mathcal{L}_{\mathrm{BPR}}=
    -
    \sum_{(u,i,j)\in\mathcal{D}}
    \log\sigma\left(
        s_{\mathrm{MF}}(u,i)-s_{\mathrm{MF}}(u,j)
    \right)
    +
    \lambda\|\Theta_{\mathrm{MF}}\|_2^2.
$$
Здесь $i$ обозначает положительный фильм пользователя $u$, а $j$ обозначает негативно
семплированный фильм, $\Theta_{\mathrm{MF}}$ обозначает множество параметров модели,
$\lambda$ обозначает коэффициент регуляризации, а $\sigma$ обозначает сигмоиду.

В данной работе матричная факторизация используется как первое базовое решение. Для
проверки влияния явных оценок также рассматривается вариант, в котором вклад
положительного взаимодействия зависит от рейтинга:
$$
    \mathcal{L}_{\mathrm{MF}}=
    -
    \sum_{(u,i,j)\in\mathcal{D}}
    \omega(r_{ui})
    \log\sigma\left(
        s_{\mathrm{MF}}(u,i)-s_{\mathrm{MF}}(u,j)
    \right)
    +
    \lambda\|\Theta_{\mathrm{MF}}\|_2^2.
$$
Если $\omega(r)=1$, то все положительные взаимодействия используются одинаково.
Если $\omega(r)$ возрастает с ростом рейтинга, то высокие оценки сильнее влияют
на обучение модели.

Дополнительно в работе проверяется использование признаков фильмов в матричной
модели. Пусть $\mathbf{x}_i\in\mathbb{R}^{f}$ обозначает вектор признаков фильма.
Первый вариант состоит в инициализации векторного представления фильма на основе признаков:
$$
    \mathbf{q}_i^{(0)}=
    W_x\mathbf{x}_i,
$$
где $W_x\in\mathbb{R}^{d\times f}$ обозначает обучаемую или заранее оценённую матрицу
проекции. Второй вариант использует признаки как добавку к обучаемому
представлению идентификатора:
$$
    \widetilde{\mathbf{q}}_i=
    \mathbf{q}_i+
    W_x\mathbf{x}_i.
$$
Тогда оценка релевантности вычисляется как
$$
    s_{\mathrm{MF+feat}}(u,i)=
    b_u+b_i+
    \left\langle
        \mathbf{p}_u,
        \widetilde{\mathbf{q}}_i
    \right\rangle.
$$
Такое расширение позволяет проверить, дают ли дополнительные признаки фильмов
прирост качества по сравнению с чисто коллаборативной моделью. Использование признаков объектов особенно важно в постановках холодного старта, что также лежит в основе гибридных моделей наподобие LightFM~\cite{kula2015lightfm}.

Матричная факторизация важна как сильное и простое базовое решение. Если более
сложная модель не превосходит её по качеству, то усложнение архитектуры не
является обоснованным.

\subsection{Двубашенные нейросетевые модели и последовательное описание истории}

Двубашенные модели являются одним из распространённых подходов для
стадии отбора кандидатов рекомендательных систем~\cite{covington2016youtube} и
продолжают линию нейросетевых моделей коллаборативной
фильтрации~\cite{he2017ncf}. Общее устройство такой модели соответствует её
названию: архитектура состоит из двух отдельных подсетей, или башен. Первая
башня кодирует пользователя и строит его векторное представление по признакам
пользователя или его истории взаимодействий. Вторая башня кодирует объект и
строит векторное представление фильма по его идентификатору, признакам или их
комбинации. После этого пользовательский и объектный векторы сравниваются в
общем латентном пространстве.

Формально одна часть модели строит представление пользователя, другая строит
представление объекта:
$$
    \mathbf{u}_{\theta}
    =
    f_{\theta}
    \left(
        \mathbf{x}^{\mathrm{user}}_u
    \right),
    \qquad
    \mathbf{v}_{\phi,i}
    =
    g_{\phi}
    \left(
        \mathbf{x}^{\mathrm{item}}_i
    \right),
$$
где
$$
    \mathbf{u}_{\theta}\in\mathbb{R}^{d},
    \qquad
    \mathbf{v}_{\phi,i}\in\mathbb{R}^{d}.
$$
Релевантность фильма $i$ для пользователя $u$ оценивается скалярным
произведением:
$$
    s_{\mathrm{TT}}(u,i)
    =
    \left\langle
        \mathbf{u}_{\theta},
        \mathbf{v}_{\phi,i}
    \right\rangle.
$$

Главное преимущество такой архитектуры состоит в масштабируемости. Представления
фильмов можно заранее вычислить офлайн, после чего для пользовательского
вектора искать ближайшие фильмы с помощью приближённого поиска ближайших
соседей. Поэтому двубашенные модели хорошо подходят для первой стадии
рекомендаций, где требуется быстро сузить большой каталог до небольшого
множества кандидатов.

В данной работе в качестве нейросетевого базового решения используется SASRec-подобная
модель~\cite{kang2018sasrec}. Такой выбор связан с тем, что последовательные
модели с механизмом внимания учитывают порядок пользовательской истории и при этом
сохраняют совместимость с оценкой релевантности через скалярное произведение.

Пусть история пользователя представлена как упорядоченная последовательность
оценённых фильмов:
$$
    \mathcal{S}_u=
    \left(
        (i_1,r_{ui_1},t_{ui_1}),
        (i_2,r_{ui_2},t_{ui_2}),
        \ldots,
        (i_L,r_{ui_L},t_{ui_L})
    \right),
$$
где
$$
    t_{ui_1}<t_{ui_2}<\ldots<t_{ui_L}.
$$
Здесь $i_l$ обозначает фильм на позиции $l$, $r_{ui_l}$ обозначает поставленный пользователем
рейтинг, а $t_{ui_l}$ обозначает время оценки.

В классическом варианте SASRec каждому объекту обычно сопоставляется обучаемое
векторное представление идентификатора. В данной работе используется другая схема: начальное представление
фильма задаётся на основе его признакового описания
$\mathbf{x}_{i}\in\mathbb{R}^{f}$:
$$
    \mathbf{v}^{(0)}_{i}
    =
    g_{\phi}
    \left(
        \mathbf{x}_{i}
    \right),
    \qquad
    \mathbf{v}^{(0)}_{i}\in\mathbb{R}^{d},
$$
где $g_{\phi}$ обозначает обучаемое отображение признаков фильма в пространство
размерности $d$. Таким образом, признаки фильма используются не как
дополнительный вектор, конкатенируемый с векторным представлением идентификатора, а как способ
инициализации векторного представления фильма.

После инициализации векторное представление фильма участвует в обучении модели и используется
как во входной последовательности пользователя, так и при оценке
фильмов-кандидатов. Для учёта порядка взаимодействий к представлению фильма на
позиции $l$ добавляется позиционное представление
$\mathbf{p}_l\in\mathbb{R}^{d}$:
$$
    \mathbf{a}_{u,l}
    =
    \mathbf{v}_{i_l}
    +
    \mathbf{p}_l.
$$

Последовательность входных векторов
$$
    A_u=
    \left(
        \mathbf{a}_{u,1},
        \mathbf{a}_{u,2},
        \ldots,
        \mathbf{a}_{u,L}
    \right)
$$
поступает в несколько блоков механизма внимания с маской. В формулах такой блок обозначается как TransformerBlock:
$$
    H^{(0)}=A_u,
$$
$$
    H^{(b+1)}
    =
    \operatorname{TransformerBlock}
    \left(
        H^{(b)}
    \right),
    \qquad
    b=0,\ldots,B-1.
$$
Маска запрещает позиции $l$ использовать информацию о будущих
элементах последовательности. Поэтому скрытое состояние $\mathbf{h}_{u,l}$ на
позиции $l$ строится только по префиксу истории:
$$
    (i_1,i_2,\ldots,i_l).
$$

Для ранжирования кандидата $i$ используется скалярное произведение скрытого
состояния пользователя и векторного представления фильма-кандидата:
$$
    s(u,l,i)
    =
    \left\langle
        \mathbf{h}_{u,l},
        \mathbf{v}_{i}
    \right\rangle.
$$
В задаче предсказания следующего объекта положительным примером для позиции
$l$ является фильм $i_{l+1}$. При негативном семплировании функция потерь может
быть записана как
$$
    \mathcal{L}_{\mathrm{SASRec}}
    =
    -
    \sum_{u}
    \sum_{l=1}^{L-1}
    \left[
        \log\sigma
        \left(
            s(u,l,i_{l+1})
        \right)
        +
        \sum_{j\in\mathcal{N}_{u,l}}
        \log
        \left(
            1-
            \sigma
            \left(
                s(u,l,j)
            \right)
        \right)
    \right],
$$
где $\mathcal{N}_{u,l}$ обозначает множество негативно семплированных фильмов для
пользователя $u$ на позиции $l$.


\begin{figure}[H]
    \centering
    \includegraphics[width=0.92\textwidth]{figures/SASRec_like.png}
    \caption{SASRec-подобная архитектура: пользовательское представление строится по последовательности оценённых фильмов с помощью блока механизма внимания. Векторное представление фильма инициализируется признаками фильма, а релевантность фильма-кандидата вычисляется через скалярное произведение пользовательского представления и представления фильма.}
    \label{fig:sasrec_two_tower_scheme}
\end{figure}

Ограничение такой модели состоит в том, что она строит представление пользователя
по его собственной последовательности и не выполняет явное распространение
информации по графу пользователь--фильм. Следовательно, часть структуры
взаимодействий, связанная с соседями более высокого порядка, может не
использоваться напрямую. Это мотивирует дальнейшее рассмотрение графовых
нейронных сетей и гибридных моделей.
\subsection{Графовые нейронные сети для рекомендаций}

Историю взаимодействий пользователей с фильмами можно представить в виде
двудольного графа
$$
    G=(\mathcal{U}\cup\mathcal{I},E),
$$
где ребро $(u,i)\in E$ означает, что пользователь $u$ оценил фильм $i$.
Графовые нейронные сети позволяют распространять информацию между соседними
узлами и учитывать не только прямые взаимодействия, но и связи более высокого
порядка. Данный подход применяется как в рекомендациях масштаба веб‑сервисов (web‑scale recommendations), так и в задачах рекомендаций в социальных сетях (social recommendation) ~\cite{ying2018pinsage,fan2019graphrec}.

Одной из базовых графовых моделей рекомендаций является NGCF~\cite{wang2019ngcf}.
В ней представление пользователя на следующем слое строится по его соседям:
$$
    \mathbf{e}^{(l+1)}_u=
    \sigma\left(
        W_1^{(l)}\mathbf{e}^{(l)}_u
        +
        \sum_{i\in N(u)}
        \frac{1}{\sqrt{|N(u)|\,|N(i)|}}
        \left(
            W_1^{(l)}\mathbf{e}^{(l)}_i
            +
            W_2^{(l)}
            \left(
                \mathbf{e}^{(l)}_i\odot\mathbf{e}^{(l)}_u
            \right)
        \right)
    \right).
$$
Здесь $N(u)$ обозначает множество соседей пользователя $u$, $\odot$ обозначает поэлементное
произведение, а $\sigma$ обозначает нелинейную функцию активации. Аналогично
обновляются представления фильмов.

\begin{figure}[H]
    \centering
    \includegraphics[width=0.87\textwidth]{figures/NGCF.png}
    \caption{Схема NGCF: векторные представления пользователя и фильма обновляются через несколько слоёв распространения сообщений, после чего представления разных слоёв объединяются и используются для предсказания релевантности пары пользователь--фильм.}
    \label{fig:ngcf_architecture}
\end{figure}

Более простой и часто более устойчивой моделью является LightGCN~\cite{he2020lightgcn}.
Она убирает нелинейные преобразования и оставляет только нормированное
распространение векторных представлений:
$$
    \mathbf{e}^{(k+1)}_u=
    \sum_{i\in N(u)}
    \frac{1}{\sqrt{|N(u)|\,|N(i)|}}
    \mathbf{e}^{(k)}_i,
$$
$$
    \mathbf{e}^{(k+1)}_i=
    \sum_{u\in N(i)}
    \frac{1}{\sqrt{|N(i)|\,|N(u)|}}
    \mathbf{e}^{(k)}_u.
$$
Финальное представление получается как сумма или среднее представлений разных
слоёв:
$$
    \mathbf{z}_u=
    \sum_{k=0}^{L}\alpha_k\mathbf{e}^{(k)}_u,
    \qquad
    \mathbf{z}_i=
    \sum_{k=0}^{L}\alpha_k\mathbf{e}^{(k)}_i.
$$
Оценка релевантности задаётся скалярным произведением:
$$
    s_{\mathrm{GNN}}(u,i)=
    \langle
        \mathbf{z}_u,
        \mathbf{z}_i
    \rangle.
$$

Помимо NGCF и LightGCN, в литературе рассматриваются и другие направления развития графовых рекомендательных моделей: разложение представлений на независимые пользовательские намерения~\cite{wang2020dgcf}, самообучение и контрастивное обучение на графах~\cite{wu2021sgl,yu2022simgcl,cai2023lightgcl}, а также подавление шума и аугментация коллаборативных графовых сигналов~\cite{fan2023graphda}. Эти методы показывают, что качество графовой рекомендации существенно зависит не только от самой схемы распространения сообщений, но и от того, как модель обрабатывает шум, разреженность и неоднородность пользовательских взаимодействий.

Графовые модели позволяют явно использовать структуру взаимодействий. Однако в
стандартной постановке ребро графа бинарное: оно показывает только факт
взаимодействия. Для датасетов с явными оценками этого недостаточно, поскольку
рейтинг несёт дополнительную информацию о силе предпочтения. Поэтому далее
рассматриваются способы включения рейтингов в графовые модели.

\begin{figure}[H]
    \centering
    \includegraphics[width=0.92\textwidth]{figures/message_pass.png}
    \caption{Идея многослойного message passing в графовой рекомендательной модели: после нескольких шагов распространения представление целевого узла агрегирует информацию от соседей более высокого порядка.}
    \label{fig:user_item_graph}
\end{figure}

\subsection{Учёт рейтингов в NGCF через весовую функцию}

В стандартной постановке NGCF граф пользователь--фильм рассматривается как
невзвешенный: ребро $(u,i)$ показывает только факт взаимодействия пользователя
$u$ с фильмом $i$. Однако в датасете с явными оценками каждому ребру
соответствует рейтинг $r_{ui}$, который содержит информацию о силе предпочтения.
Поскольку рейтинги образуют упорядоченную шкалу, их можно использовать не только как независимые категории, но и как порядковый сигнал. Такой подход соответствует идее учёта порядковости рейтингов в графовых моделях матричного восстановления~\cite{lee2024rogmc}. Поэтому в данной работе рассматривается модификация NGCF, в которой рейтинг
преобразуется в вес ребра:
$$
    w_{ui}=\omega(r_{ui}).
$$
Тогда фильмы, получившие более высокие оценки, сильнее влияют на представление
пользователя при распространении сообщений.

В стандартном NGCF сообщение от фильма $i$ к пользователю $u$ имеет вид
$$
    \mathbf{m}^{(l)}_{u\leftarrow i}
    =
    \frac{1}{\sqrt{|N(u)|\,|N(i)|}}
    \left(
        W_1^{(l)}\mathbf{e}^{(l)}_i
        +
        W_2^{(l)}
        \left(
            \mathbf{e}^{(l)}_i\odot \mathbf{e}^{(l)}_u
        \right)
    \right),
$$
где $N(u)$ и $N(i)$ обозначают множества соседей пользователя и фильма,
$\mathbf{e}^{(l)}_u,\mathbf{e}^{(l)}_i\in\mathbb{R}^{d_l}$ обозначают их
представления на слое $l$, $W_1^{(l)}$ и $W_2^{(l)}$ являются обучаемыми матрицами, а
$\odot$ обозначает поэлементное произведение.

В рейтинговой модификации нормировочный коэффициент заменяется на
взвешенный:
$$
    \widetilde A_{ui}
    =
    \frac{
        \omega(r_{ui})
    }{
        \sqrt{
            \sum\limits_{i'\in N(u)}
            \omega(r_{ui'})
        }
        \sqrt{
            \sum\limits_{u'\in N(i)}
            \omega(r_{u'i})
        }
    }.
$$
Тогда сообщение принимает вид
$$
    \mathbf{m}^{(l)}_{u\leftarrow i}
    =
    \widetilde A_{ui}
    \left(
        W_1^{(l)}\mathbf{e}^{(l)}_i
        +
        W_2^{(l)}
        \left(
            \mathbf{e}^{(l)}_i\odot \mathbf{e}^{(l)}_u
        \right)
    \right).
$$
После агрегации сообщений обновление пользователя записывается как
$$
    \mathbf{e}^{(l+1)}_u
    =
    \sigma
    \left(
        W_1^{(l)}\mathbf{e}^{(l)}_u
        +
        \sum_{i\in N(u)}
        \mathbf{m}^{(l)}_{u\leftarrow i}
    \right).
$$
Аналогично строится обновление представления фильма $i$ по соседним
пользователям. Таким образом, архитектура NGCF сохраняется, но интенсивность
передачи информации по ребру зависит от явного рейтинга.

В работе рассматриваются три способа задания функции $\omega(r)$:
линейная функция, экспоненциальная функция и функция на основе softplus.

\paragraph{Линейная функция рейтингового веса.}

Первый вариант использует нормированную линейную функцию:
$$
    \omega_{\mathrm{lin}}(r)
    =
    \frac{r-r_{\min}}{r_{\max}-r_{\min}}.
$$
Здесь $r_{\min}$ и $r_{\max}$ обозначают минимальное и максимальное возможные значения
рейтинга. Для шкалы MovieLens~\cite{harper2015movielens} можно считать $r_{\min}=0.5$ и $r_{\max}=5.0$.
В этом случае минимальная оценка получает вес, близкий к нулю, а максимальная
оценка получает вес, равный единице.

Преимущество линейной функции состоит в простоте и отсутствии дополнительных
обучаемых параметров. Она сохраняет порядок рейтингов и задаёт естественную
интерпретацию: чем выше оценка, тем сильнее соответствующее ребро в графе.
Недостаток такого подхода состоит в том, что форма зависимости фиксирована
заранее. Например, модель не может сама определить, что разница между оценками
$4.5$ и $5.0$ важнее или менее важна, чем разница между $2.0$ и $2.5$.

\paragraph{Экспоненциальная функция рейтингового веса.}

Второй вариант использует экспоненциальное преобразование рейтинга:
$$
    \omega_{\mathrm{exp}}(r)
    =
    \exp\left(
        \alpha
        \frac{r-r_{\min}}{r_{\max}-r_{\min}}
    \right),
$$
где $\alpha$ обозначает параметр, управляющий резкостью зависимости. Чтобы средний
масштаб весов оставался сопоставимым с другими вариантами, веса можно
нормировать:
$$
    \bar\omega_{\mathrm{exp}}(r)
    =
    \frac{
        \omega_{\mathrm{exp}}(r)
    }{
        \frac{1}{M}
        \sum\limits_{m=1}^{M}
        \omega_{\mathrm{exp}}(r_m)
    },
$$
где $r_1,\ldots,r_M$ обозначают все возможные значения рейтинга.

Экспоненциальная функция усиливает различия между высокими оценками. При
$\alpha>0$ рост веса становится нелинейным: высокие рейтинги получают
существенно больший вклад в распространение сообщений, чем средние и низкие.
Такой вариант проверяет гипотезу о том, что для построения списка из $K$ рекомендаций особенно
важны не все положительные оценки, а наиболее сильные предпочтения
пользователя.

Если параметр $\alpha$ фиксируется заранее, то функция остаётся простой
эвристикой. Если $\alpha$ обучается, то модель получает возможность подобрать
степень усиления высоких рейтингов по данным.

\paragraph{Softplus-функция рейтингового веса.}

Третий вариант задаёт вес через softplus-преобразование:
$$
    \omega_{\mathrm{sp}}(r)
    =
    \operatorname{softplus}
    \left(
        a r + b
    \right)
    +
    \varepsilon,
$$
где $a$ и $b$ являются обучаемыми параметрами, а $\varepsilon>0$ добавляется для
численной устойчивости. Напомним, что
$$
    \operatorname{softplus}(x)
    =
    \log(1+\exp(x)).
$$
Эта функция положительна и гладко аппроксимирует $\max(x,0)$. Поэтому она
удобна для задания весов рёбер: вес не становится отрицательным, а зависимость
от рейтинга может подстраиваться под данные.

Для сохранения порядка рейтингов можно ограничить параметр $a$ быть
неотрицательным:
$$
    a=\operatorname{softplus}(\eta),
$$
где $\eta$ является обучаемым параметром. Тогда
$$
    r_a\leq r_b
    \quad\Rightarrow\quad
    \omega_{\mathrm{sp}}(r_a)
    \leq
    \omega_{\mathrm{sp}}(r_b).
$$
Иными словами, больший рейтинг не может получить меньший вес. Как и для
экспоненциального варианта, веса можно нормировать:
$$
    \bar\omega_{\mathrm{sp}}(r)
    =
    \frac{
        \omega_{\mathrm{sp}}(r)
    }{
        \frac{1}{M}
        \sum\limits_{m=1}^{M}
        \omega_{\mathrm{sp}}(r_m)
    }.
$$

Softplus-функция является более гибкой, чем линейная: она позволяет модели
выучить как почти линейную зависимость, так и зависимость с пороговым эффектом,
при которой низкие рейтинги почти не влияют на распространение сообщений, а
высокие получают заметно больший вес.

\paragraph{Итоговая схема рейтингового NGCF.}

Все три варианта отличаются только способом вычисления веса $\omega(r_{ui})$.
Дальнейшая архитектура остаётся одинаковой: по рейтингам строится взвешенная
нормировка $\widetilde A_{ui}$, после чего она используется в NGCF-сообщениях.
Финальные представления пользователя и фильма после $L$ слоёв обозначаются как
$\mathbf{z}_u$ и $\mathbf{z}_i$, а оценка релевантности вычисляется скалярным произведением:
$$
    s_{\mathrm{NGCF}}(u,i)
    =
    \left\langle
        \mathbf{z}_u,
        \mathbf{z}_i
    \right\rangle.
$$
Обучение выполняется с помощью BPR-функции потерь:
$$
    \mathcal{L}_{\mathrm{NGCF}}
    =
    -
    \sum_{(u,i,j)\in\mathcal{D}}
    \log\sigma
    \left(
        s_{\mathrm{NGCF}}(u,i)
        -
        s_{\mathrm{NGCF}}(u,j)
    \right)
    +
    \lambda
    \|\Theta_{\mathrm{NGCF}}\|_2^2.
$$
Здесь $i$ обозначает положительный фильм пользователя, $j$ обозначает негативно
семплированный фильм, а $\Theta_{\mathrm{NGCF}}$ обозначает параметры модели.

Такое сравнение позволяет отделить влияние самой графовой архитектуры от
влияния способа преобразования рейтинга в вес ребра. Линейная функция служит
простой базовой схемой, экспоненциальная функция проверяет пользу усиления
высоких оценок, а функция на основе softplus позволяет обучать форму зависимости по
данным при сохранении положительности весов.

\subsection{GCMC: рейтинг как тип ребра}

Другой способ учёта рейтинга состоит в том, чтобы рассматривать каждое значение
рейтинга как отдельный тип ребра. Такая идея используется в модели Graph
Convolutional Matrix Completion~\cite{vandenberg2017gcmc}. Пусть
$\mathcal{R}$ обозначает множество возможных рейтингов. Для каждого рейтинга
$r\in\mathcal{R}$ строится отдельная матрица смежности $M_r$. Ребро типа $r$
соединяет пользователя $u$ и фильм $i$, если пользователь поставил этому
фильму оценку $r$.

Для каждого типа рейтинга используется своя матрица параметров $W_r$.
Сообщение от фильма $i$ к пользователю $u$ имеет вид
$$
    \boldsymbol{\mu}_{i\to u,r}=
    \frac{1}{c_{ui}}
    W_r\mathbf{x}_i,
$$
где $c_{ui}$ обозначает нормировочный коэффициент. После агрегации сообщений разных
типов получаем скрытое представление пользователя:
$$
    \mathbf{h}_u=
    \sigma\left(
        \operatorname{accum}_{r\in\mathcal{R}}
        \sum_{i\in N_r(u)}
        \boldsymbol{\mu}_{i\to u,r}
    \right),
$$
$$
    \mathbf{z}_u=
    \sigma(W\mathbf{h}_u).
$$
Аналогично строится представление фильма $\mathbf{z}_i$.

Decoder предсказывает распределение по возможным рейтингам:
$$
    p_{\theta}(\hat r_{ui}=r)=
    \frac{
        \exp\left(
            \mathbf{z}_u^{\top}
            Q_r
            \mathbf{z}_i
        \right)
    }{
        \sum\limits_{s\in\mathcal{R}}
        \exp\left(
            \mathbf{z}_u^{\top}
            Q_s
            \mathbf{z}_i
        \right)
    },
$$
где $Q_r\in\mathbb{R}^{d\times d}$ обозначает матрицу декодировщика для рейтинга $r$.
Ожидаемый рейтинг:
$$
    \hat r_{ui}=
    \sum_{r\in\mathcal{R}}
    r\,
    p_{\theta}(\hat r_{ui}=r).
$$

Изначально GCMC обучается на задачу предсказания рейтинга. Для этого
используется функция потерь перекрёстной энтропии по наблюдаемым оценкам:
$$
    \mathcal{L}_{\mathrm{CE}}=
    -
    \sum_{(u,i)\in\Omega}
    \sum_{r\in\mathcal{R}}
    \mathbb{I}[r_{ui}=r]
    \log p_{\theta}(\hat r_{ui}=r).
$$
Однако целевая задача данной работы состоит не в точном предсказании численного
рейтинга, а в ранжировании фильмов для построения первых $K$ позиций списка. Поэтому в реализованной модификации GCMC
задача предсказания рейтинга дополняется ранжирующей компонентой:
$$
    \mathcal{L}_{\mathrm{GCMC}}=
    \mathcal{L}_{\mathrm{CE}}
    +
    \lambda_{\mathrm{rank}}
    \mathcal{L}_{\mathrm{BPR}}.
$$
Здесь $\lambda_{\mathrm{rank}}$ управляет вкладом ранжирующей функции потерь и
подбирается по качеству модели на валидационной выборке. Идея совместного
использования задачи ранжирования и задачи предсказания рейтинга близка к
многозадачной постановке Rank and Rate~\cite{hadash2018rankrate}. Такая
модификация позволяет сохранить способность GCMC учитывать типы рейтинговых
рёбер и одновременно оптимизировать порядок объектов в рекомендательном списке.
Конкретный баланс между двумя целями выбирается на валидации.

\paragraph{Временная модификация GCMC.}

Также будем рассматривать модификацию GCMC  с помощью учёта времени. Для
взаимодействия $(u,i,r_{ui},t_{ui})$ вводится коэффициент затухания:
$$
    \alpha(t,t_{ui})=
    \exp\left(
        -\gamma(t-t_{ui})
    \right),
$$
где $t$ обозначает момент построения рекомендации, а $\gamma>0$ обозначает параметр скорости
затухания. Тогда сообщение в GCMC принимает вид
$$
    \boldsymbol{\mu}_{i\to u,r}(t)=
    \frac{
        \alpha(t,t_{ui})
    }{
        c_{ui}(t)
    }
    W_r\mathbf{x}_i.
$$
Если дополнительно используется вес рейтинга, то коэффициент можно записать как
$$
    a_{ui}(t)=
    \omega(r_{ui})
    \alpha(t,t_{ui}).
$$
Тогда сообщение имеет вид
$$
    \boldsymbol{\mu}_{i\to u,r}(t)=
    \frac{
        a_{ui}(t)
    }{
        c_{ui}(t)
    }
    W_r\mathbf{x}_i.
$$
Такая модификация отражает предположение, что более свежие оценки лучше
характеризуют текущие предпочтения пользователя, чем старые.

\begin{figure}[H]
    \centering
    \includegraphics[width=0.92\textwidth]{figures/GCMC.png}
    \caption{Архитектура GCMC: графовый кодировщик строит пользовательские и представления фильма с учётом рейтинговых связей, а bilinear декодировщик восстанавливает матрицу оценок.}
    \label{fig:gcmc_scheme}
\end{figure}

\subsection{Гибридные модели: совместное использование двубашенных и графовых представлений}

Двубашенная и графовая архитектуры используют разные механизмы построения
представлений. Двубашенная модель хорошо масштабируется и позволяет независимо
кодировать пользователя и фильм. Графовая модель, напротив, строит
представления через распространение информации по графу взаимодействий
пользователь--фильм. Поэтому в работе рассматриваются гибридные подходы, в
которых выход одной модели используется как вход другой модели.

В отличие от простого добавления признаков, в этих вариантах представления
одной архитектуры подаются непосредственно на вход другой. Таким образом,
проверяется не только возможность объединения оценок релевантности, но и то, может ли одна
модель улучшить входное представление для другой.

\paragraph{Графовые представления как вход двубашенной модели.}

В первом варианте сначала строятся графовые представления пользователя и фильма:
$$
    \mathbf{z}^{\mathrm{GNN}}_u,
    \qquad
    \mathbf{z}^{\mathrm{GNN}}_i.
$$
После этого они используются как входы пользовательской башни и башни фильма:
$$
    \mathbf{u}^{\mathrm{G\to TT}}_u=
    f_{\theta}
    \left(
        \mathbf{z}^{\mathrm{GNN}}_u
    \right),
$$
$$
    \mathbf{v}^{\mathrm{G\to TT}}_i=
    g_{\phi}
    \left(
        \mathbf{z}^{\mathrm{GNN}}_i
    \right).
$$
Оценка релевантности вычисляется как скалярное произведение:
$$
    s_{\mathrm{G\to TT}}(u,i)=
    \left\langle
        \mathbf{u}^{\mathrm{G\to TT}}_u,
        \mathbf{v}^{\mathrm{G\to TT}}_i
    \right\rangle.
$$

Смысл этого варианта состоит в том, что графовая модель предварительно
извлекает информацию из структуры взаимодействий, а двубашенная модель затем
переотображает полученные графовые представления в общее пространство
ранжирования. При этом двубашенная архитектура не получает исходные
признаки идентификаторов или признаки фильмов напрямую: её входом являются именно
представления, полученные графовой моделью.

\paragraph{Двубашенные представления как вход графовой модели.}

Во втором варианте сначала строятся представления пользователя и фильма с
помощью двубашенной модели:
$$
    \mathbf{u}^{\mathrm{TT}}_u,
    \qquad
    \mathbf{v}^{\mathrm{TT}}_i.
$$
Затем эти векторы используются как начальные представления узлов графовой
модели:
$$
    \mathbf{e}^{(0)}_u=
    \mathbf{u}^{\mathrm{TT}}_u,
    \qquad
    \mathbf{e}^{(0)}_i=
    \mathbf{v}^{\mathrm{TT}}_i.
$$
После этого выполняется графовое распространение. В NGCF-подобной модели
обновление пользователя имеет вид
$$
    \mathbf{e}^{(l+1)}_u=
    \sigma\left(
        W_1^{(l)}\mathbf{e}^{(l)}_u
        +
        \sum_{i\in N(u)}
        \widetilde A_{ui}
        \left(
            W_1^{(l)}\mathbf{e}^{(l)}_i
            +
            W_2^{(l)}
            \left(
                \mathbf{e}^{(l)}_i\odot\mathbf{e}^{(l)}_u
            \right)
        \right)
    \right).
$$
Аналогично строится обновление представления фильма по соседним пользователям.
Финальные графовые представления обозначим как
$$
    \mathbf{z}^{\mathrm{TT\to G}}_u,
    \qquad
    \mathbf{z}^{\mathrm{TT\to G}}_i.
$$
Оценка релевантности имеет вид
$$
    s_{\mathrm{TT\to G}}(u,i)=
    \left\langle
        \mathbf{z}^{\mathrm{TT\to G}}_u,
        \mathbf{z}^{\mathrm{TT\to G}}_i
    \right\rangle.
$$

В этом варианте двубашенная модель сначала формирует входные представления
пользователей и фильмов, а графовая модель затем распространяет их по графу
взаимодействий. То есть признаки и история пользователя учитываются до
графового этапа, а последующее распространение позволяет уточнить полученные
векторы с учётом соседей в графе.

\paragraph{Взвешенная сумма оценок релевантности и представлений.}

Третий вариант не передаёт представления одной модели внутрь другой, а
объединяет уже полученные оценки релевантности:
$$
    s_{\mathrm{ens}}(u,i)=
    \beta s_{\mathrm{TT}}(u,i)
    +
    (1-\beta)s_{\mathrm{GNN}}(u,i),
$$
где $\beta\in[0,1]$ обозначает коэффициент смешивания. Он подбирается по
валидационной выборке:
$$
    \beta^{*}=
    \arg\max_{\beta\in[0,1]}
    \operatorname{Metric}_{\mathrm{val}}
    \left(
        s_{\mathrm{ens}}
    \right).
$$

В экспериментах также используется близкий вариант ансамблирования на уровне
представлений. Пусть двубашенная и графовая модели строят векторы
$\mathbf{u}^{\mathrm{TT}}_u$, $\mathbf{v}^{\mathrm{TT}}_i$ и
$\mathbf{u}^{\mathrm{GNN}}_u$, $\mathbf{v}^{\mathrm{GNN}}_i$. Сначала они
проецируются в общее пространство размерности $d$:
$$
    \widetilde{\mathbf{u}}^{\mathrm{TT}}_u=P^{\mathrm{TT}}_u
    \mathbf{u}^{\mathrm{TT}}_u,
    \qquad
    \widetilde{\mathbf{u}}^{\mathrm{GNN}}_u=P^{\mathrm{GNN}}_u
    \mathbf{u}^{\mathrm{GNN}}_u,
$$
$$
    \widetilde{\mathbf{v}}^{\mathrm{TT}}_i=P^{\mathrm{TT}}_i
    \mathbf{v}^{\mathrm{TT}}_i,
    \qquad
    \widetilde{\mathbf{v}}^{\mathrm{GNN}}_i=P^{\mathrm{GNN}}_i
    \mathbf{v}^{\mathrm{GNN}}_i.
$$
После этого пользовательское и представления фильма смешиваются с отдельными
коэффициентами:
$$
    \mathbf{u}^{\mathrm{ens}}_u=
    \alpha_u\widetilde{\mathbf{u}}^{\mathrm{TT}}_u+
    (1-\alpha_u)\widetilde{\mathbf{u}}^{\mathrm{GNN}}_u,
$$
$$
    \mathbf{v}^{\mathrm{ens}}_i=
    \alpha_i\widetilde{\mathbf{v}}^{\mathrm{TT}}_i+
    (1-\alpha_i)\widetilde{\mathbf{v}}^{\mathrm{GNN}}_i,
$$
где $\alpha_u,\alpha_i\in[0,1]$. На практике коэффициенты задаются через
обучаемые скаляры
$$
    \alpha_u=\mathrm{sigmoid}(a_u),
    \qquad
    \alpha_i=\mathrm{sigmoid}(a_i),
$$
что удерживает их в допустимом диапазоне. Финальная оценка релевантности имеет вид
$$
    s_{\mathrm{ens}}(u,i)=
    \left\langle
        \mathbf{u}^{\mathrm{ens}}_u,
        \mathbf{v}^{\mathrm{ens}}_i
    \right\rangle.
$$
Если $\alpha_u$ или $\alpha_i$ близки к $1$, ансамбль сильнее опирается на
двубашенную компоненту. Если они близки к $0$, больший вклад получает графовая
компонента. Значения около $0.5$ соответствуют почти равному смешиванию двух
источников представлений.

Этот вариант является простым способом проверить, дают ли двубашенная и
графовая архитектуры взаимодополняющие оценки релевантности. В отличие от
полного каскадного гибрида, он не заставляет одну модель полностью
перестраивать представления другой, а объединяет два уже обученных источника
сигнала.

Для честного сравнения гибридные модели оцениваются не только по качеству, но
и по времени обучения и применения. Это важно, поскольку объединение нескольких
моделей может повышать HitRate@$K$ и nDCG@$K$, но одновременно увеличивать
вычислительные затраты.

\subsection{Сводная таблица рассматриваемых моделей}

\begin{table}[H]
\centering
\caption{Классы моделей, рассматриваемые в работе.}
\label{tab:models_overview}
\begin{tabular}{|p{0.22\textwidth}|p{0.34\textwidth}|p{0.33\textwidth}|}
\hline
\textbf{Модель} &
\textbf{Используемый сигнал} &
\textbf{Назначение в экспериментах} \\
\hline
MF &
Вектор пользователя и фильма, BPR-ранжирование &
Базовая модель коллаборативной фильтрации. \\
\hline
MF с признаками &
Вектор пользователя, вектор фильма, признаки фильма &
Проверка пользы признаков фильмов в простой модели. \\
\hline
SASRec &
Последовательность оценённых фильмов пользователя, позиционные представления, механизм внимания &
Сильное нейросетевое последовательное базовое решение, близкое по идее к моделям отбора кандидатов для персонализированного предложения. \\
\hline
NGCF &
Граф пользователь--фильм &
Базовые графовые модели рекомендаций. \\
\hline
NGCF с рейтинговыми весами &
Рейтинг как вес ребра &
Проверка линейной, экспоненциальной и softplus-функций рейтингового веса. \\
\hline
GCMC &
Рейтинг как тип ребра &
Проверка способа учёта явных оценок через типы рейтинговых связей. \\
\hline
GCMC с учётом времени &
Тип рейтингового ребра и временное затухание &
Проверка пользы учёта времени в GCMC. \\
\hline
Гибридные модели &
Передача выходов графовой нейронной сети на вход двубашенной модели, передача выходов двубашенной модели на вход графовой нейронной сети,
взвешенная сумма оценок релевантности &
Проверка взаимодополняемости двубашенных и графовых моделей. \\
\hline
\end{tabular}
\end{table}



\section{Эксперименты и результаты}
\label{sec:experiments}

В данном разделе проводится экспериментальное сравнение моделей, описанных
в таблице ~\ref{tab:models_overview}. 
Основное внимание в разделе уделяется исследованию способов учёта
рейтинга в графовых моделях: сравнению подхода с типами рейтинговых связей, где разные
значения рейтинга задают разные типы связей, порядкового подхода, в котором рейтинг
определяет силу ребра, и временной модификации, где более свежие оценки имеют
больший вклад. Дополнительно проверяется роль признаков фильмов, а также
гипотеза о том, что графовые представления могут дополнять двубашенную
модель отбора кандидатов в составе гибридов и ансамблей.


\subsection{Датасет и разбиение данных}
\label{subsec:experiment_dataset}

Эксперименты проводились на подготовленном датасете MovieLens с явными
пользовательскими оценками. Каждое взаимодействие содержит идентификатор
пользователя, идентификатор фильма, оценку и временную метку. В качестве
положительными объектами для бинарных метрик считались фильмы с
оценкой не ниже $4.0$, как определено в подразделе~\ref{subsec:metrics}.

Разбиение данных н а обучающую, валидационную и тестовую части можно описать следующим
образом:
\begin{itemize}
    \item разбиение выполнялось по времени отдельно для каждого пользователя.
    Все взаимодействия пользователя сортировались по полю \texttt{timestamp}.
    Положительной оценкой при построении разбиения считалась оценка
    $r_{ui} \geq 4$;

    \item для каждого пользователя последние $3$ положительные оценки
    относились к тестовой части, а предыдущие $2$ положительные оценки
    относились к валидационной части. Обучающая часть содержала взаимодействия,
    произошедшие раньше валидационной части. При этом в валидационную и
    тестовую части сохранялись все оценки из соответствующих временных
    промежутков, включая низкие и нейтральные;

    \item для моделей, использующих историю пользователя, брались последние
    $50$ взаимодействий пользователя перед прогнозируемым событием. Если у
    пользователя было меньше $50$ взаимодействий, использовалась вся доступная
    история;

    \item обучающая выборка тёплого старта содержит $19\,574\,307$
    взаимодействий, $144\,009$ пользователей и $31\,999$ фильмов. Валидационная
    и тестовая части для тёплого старта содержат соответственно $604\,994$ и
    $987\,722$ взаимодействия;

    \item для сценария холодного старта для клиента использовалось $15\,935$
    новых пользователей. В этом режиме пользователь не имеет обучаемого
    идентификаторного представления в основной обучающей выборке, поэтому его
    представление строится по опорной истории;

    \item опорная история для сценария холодного старта для клиента содержит
    $2\,295\,457$ взаимодействий. Валидационная часть содержит $66\,243$
    взаимодействия, а тестовая часть содержит $106\,919$ взаимодействий.
\end{itemize}

Для фильмов использовались $160$ дополнительных признаковых колонок: жанры MovieLens и IMDb,
год выпуска, десятилетие, количество жанров, SVD-сжатие Genome-тегов, а также
агрегированные признаки по актёрам, актрисам, режиссёрам и сценаристам.
Масштабирование признаков выполнялось только по части тёплого старта, чтобы не
допускать утечки информации из данных холодного старта по пользователю и тестовых данных.

Подробное описание схемы разбиения, а также собранных дополнительных признаков
фильмов приведено в приложении~\ref{app:dataset_details}.

\subsection{Протокол обучения и оценки}
\label{subsec:experiment_protocol}

Все модели сравнивались в единой постановке. При оценке
из множества кандидатов исключались фильмы, уже встречавшиеся в истории
пользователя. Основные метрики качества определены в
подразделе~\ref{subsec:metrics}. В основной части экспериментов используются
HR@20, HR@50, nDCG@20 и nDCG@50.

Обучение выполнялось с фиксированным начальным зерном генератора случайных чисел $42$. Для нейросетевых
моделей использовался метод оптимизации Adam с коэффициентом $L_2$-регуляризации $10^{-5}$. В базовых запусках
применялись скорость обучения $3\cdot 10^{-4}$, размер батча $1024$, один негативный
пример на положительное взаимодействие. Для совместного дообучения ансамбля TT и GCMC использовалась
скорость обучения $10^{-4}$.

Гиперпараметры и лучшая контрольная точка модели выбирались по валидационной выборке.
Тестовая выборка использовалась только для итогового сравнения моделей, чтобы
не допускать подбора параметров по тестовым данным. Более подробные настройки
обучения и оценки приведены в приложении~\ref{app:experiment_extra}.

\subsection{Базовые модели и влияние признаков}
\label{subsec:experiment_features}

Цель первого эксперимента состоит в том, чтобы зафиксировать простые
бейзлайны и проверить, как признаки фильмов влияют на качество моделей. Для
этого сравнивались варианты матричной факторизации с разной размерностью
эмбеддингов, а также два способа использования признаков: инициализация
векторных представлений фильмов по признаковому представлению и прямое добавление
признакового представления к латентному вектору. Дополнительно проверялись
варианты GCMC без признаков, с признаковой инициализацией и с дополнительным
добавлением признаков.

\begin{table}[htbp]
	\centering
	\footnotesize
	\setlength{\tabcolsep}{4pt}
	\caption{Влияние размерности эмбеддингов и способа использования признаков фильмов в MF. Все метрики приведены для warm-start сценария.}
	\label{tab:features_mf}
	\begin{tabular}{@{}lccc@{}}
		\toprule
		Модель & HR@20 & HR@50 & nDCG@20 \\
		\midrule
		\texttt{MF-64}      & 0.2447 & 0.3956 & 0.0601 \\
		\texttt{MF-128}     & 0.2458 & 0.3990 & 0.0616 \\
		\texttt{MF-128-init}& \textbf{0.2479} & \textbf{0.4021} & \textbf{0.0621} \\
		\texttt{MF-128-add} & 0.2349 & 0.3971 & 0.0560 \\
		\texttt{MF-256-add} & 0.2385 & 0.3969 & 0.0577 \\
		\bottomrule
	\end{tabular}
\end{table}

Из таблицы~\ref{tab:features_mf} видно, что среди MF-моделей лучший результат
показывает конфигурация \texttt{MF-128-init}: она достигает $0.2479$ по
HR@20, $0.4021$ по HR@50 и $0.0621$ по nDCG@20. По сравнению с
\texttt{MF-128} использование признаков для инициализации даёт небольшой, но
устойчивый прирост по всем трём метрикам. Это означает, что контентные
признаки фильмов помогают модели стартовать из более информативного начального
представления и быстрее сходиться к хорошему решению.

При этом прямое добавление признакового представления оказывается менее
удачным. Модели \texttt{MF-128-add} и \texttt{MF-256-add} уступают не только
варианту с признаковой инициализацией, но и обычным MF-моделям. Следовательно,
признаки полезны не сами по себе, а при аккуратном способе включения в модель:
в данной постановке они лучше работают как источник начальной структуры
векторных представлений фильмов, чем как дополнительная компонента финальной оценки релевантности.

\begin{table}[htbp]
\centering
\small
\caption{Влияние признаков фильмов в GCMC в warm-start сценарии. Эксперимент проверяет, улучшает ли прямое добавление item features качество графового распространения.}
\label{tab:features_gcmc}
\begin{tabular}{lrrr}
\toprule
Модель & Warm HR@20 & Warm HR@50 & Warm nDCG@20 \\
\midrule
\texttt{\detokenize{GCMC-noFeat}} & 0.1887 & 0.3103 & 0.0473 \\
\texttt{\detokenize{GCMC}} & \textbf{0.1889} & \textbf{0.3121} & \textbf{0.0477} \\
\texttt{\detokenize{GCMC-add}} & 0.1871 & 0.3098 & 0.0457 \\
\bottomrule
\end{tabular}
\end{table}


Похожая картина наблюдается и для GCMC. Таблица~\ref{tab:features_gcmc}
показывает, что базовая \texttt{GCMC} немного превосходит \texttt{GCMC-noFeat}
и \texttt{GCMC-add}: она достигает $0.1889$ по HR@20 и $0.0477$ по nDCG@20.
Вариант \texttt{GCMC-add}, где признаки добавлялись дополнительно, оказывается
хуже. Это подтверждает, что механическое добавление side information может
вносить шум и ухудшать согласованность графового распространения с
коллаборативным сигналом.

Вывод эксперимента состоит в том, что признаки фильмов могут быть полезны,
но эффект зависит от способа их использования. В MF наилучший результат даёт
признаковая инициализация, поскольку она задаёт более удачную стартовую точку
для векторных представлений фильмов. Однако прямое добавление признаков к латентным
представлениям в MF и GCMC не приводит к улучшению. Поэтому дальнейшие
эксперименты сосредоточены не только на признаках, но и на способах
использования явной обратной связи в графе и на объединении разных архитектур.

\subsection{Влияние временного затухания в GCMC}
\label{subsec:experiment_temporal_gcmc}

Отдельный эксперимент посвящён проверке временного сигнала в GCMC. Его цель
состоит в том, чтобы проверить гипотезу о большей важности свежих оценок:
если предпочтения пользователя меняются со временем, то более поздние
взаимодействия должны сильнее влиять на пользовательское представление.
Сравнивались базовая \texttt{GCMC} и \texttt{GCMC-time}, в которой сообщения
по рёбрам дополнительно взвешивались коэффициентом временного затухания.

\begin{table}[htbp]
\centering
\small
\caption{Влияние временного затухания.}
\label{tab:temporal_decay}
\begin{tabular}{llrrrrrr}
\toprule
Модель & Temporal decay & Warm HR@20 & Warm HR@50 & Warm nDCG@20 & Cold-user HR@20 & Cold-user HR@50 & Cold-user nDCG@20 \\
\midrule
\texttt{\detokenize{GCMC}} & нет & 0.1889 & 0.3121 & 0.0477 & 0.1892 & 0.3062 & \textbf{0.0470} \\
\texttt{\detokenize{GCMC-time}} & да & \textbf{0.1989} & \textbf{0.3280} & \textbf{0.0482} & \textbf{0.1946} & \textbf{0.3134} & 0.0469 \\
\bottomrule
\end{tabular}
\end{table}


Таблица~\ref{tab:temporal_decay} показывает, что добавление временного
затухания улучшает качество по основным HR-метрикам. HR@20 в режиме тёплого старта возрастает
с $0.1889$ до $0.1989$, то есть примерно на $5.3\%$, а HR@50 в режиме тёплого старта увеличивается с
$0.3121$ до $0.3280$. В сценарии холодного старта по пользователю HR@20 также возрастает: с
$0.1892$ до $0.1946$. При этом nDCG@20 в сценарии холодного старта по пользователю почти не меняется
($0.0470$ против $0.0469$), поэтому эффект времени проявляется прежде всего
в способности модели чаще включать релевантные фильмы в первые $K$ позиций, а не
обязательно ставить их заметно выше внутри списка.

Вывод эксперимента состоит в том, что временная информация действительно
полезна для GCMC. Более свежие оценки оказываются более информативными для
построения пользовательского представления, поэтому \texttt{GCMC-time}
используется далее как лучший GCMC-вариант в итоговом сравнении.

\subsection{Влияние функции рейтингового веса}
\label{subsec:experiment_rating_weight}

Следующий эксперимент является детализацией порядкового подхода к учёту рейтинга.
Его цель состоит в том, чтобы проверить, достаточно ли простой заранее заданной монотонной
функции веса или полезно обучать форму зависимости веса ребра от значения
рейтинга. Сравнивались три функции рейтингового веса: линейная,
экспоненциальная смесь и обучаемая функция на основе softplus-сплайна.

\begin{table}[htbp]
\centering
\small
\caption{Влияние функции рейтингового веса в NGCF в warm-start сценарии. Эксперимент проверяет, полезно ли обучать зависимость веса ребра от значения рейтинга.}
\label{tab:rating_weight_function}
\begin{tabular}{llrrr}
\toprule
Функция веса & Модель & Warm HR@20 & Warm HR@50 & Warm nDCG@20 \\
\midrule
Линейная & \texttt{\detokenize{NGCF-lin}} & 0.2024 & 0.3545 & 0.0452 \\
Экспоненциальная смесь & \texttt{\detokenize{NGCF-exp}} & 0.2038 & 0.3555 & 0.0455 \\
Обучаемая softplus-spline & \texttt{\detokenize{NGCF-spline}} & \textbf{0.2108} & \textbf{0.3644} & \textbf{0.0473} \\
\bottomrule
\end{tabular}
\end{table}


\begin{figure}[htbp]
    \centering
    \includegraphics[width=0.72\textwidth]{figures/rating_weight_learned_function}
    \caption{Обучаемая функция рейтингового веса для NGCF-модели. Функция
    сильнее выделяет высокие оценки и подавляет вклад низких оценок.}
    \label{fig:rating_weight_learned_function}
\end{figure}

Таблица~\ref{tab:rating_weight_function} показывает, что обучаемая
функция на основе softplus-сплайна даёт лучший результат в сценарии тёплого старта:
\texttt{NGCF-spline} достигает $0.2108$ по HR@20 в режиме тёплого старта, $0.3644$ по HR@50 в режиме тёплого старта
и $0.0473$ по nDCG@20 в режиме тёплого старта. Это говорит о том, что модели полезно не просто
монотонно увеличивать вес ребра с ростом рейтинга, а подбирать форму этой
зависимости по данным.

Вид обученной функции на рис.~\ref{fig:rating_weight_learned_function}
согласуется с численными результатами: модель усиливает вклад высоких оценок
и ослабляет вклад низких. Таким образом, явный рейтинг оказывается полезным
не только как бинарный индикатор положительного взаимодействия, но и как
порядковый сигнал, задающий силу связи в графе.

Вывод эксперимента состоит в том, что обучаемая функция рейтингового веса
улучшает качество NGCF-модели в режиме тёплого старта по сравнению с фиксированными
функциями. Поэтому \texttt{NGCF-spline} далее рассматривается как лучший
NGCF-вариант при общем сравнении графовых моделей. При этом устойчивость этой
модификации к новым пользователям требует отдельной проверки и анализируется
ниже в эксперименте с холодным стартом по пользователю.

\subsection{Сравнение графовых моделей}
\label{subsec:experiment_graph_models}

После отдельных проверок временного затухания в GCMC и функции рейтингового
веса в NGCF проводится общее сравнение графовых моделей. Цель эксперимента состоит в том, чтобы
сравнить два способа использования явных рейтингов в графе.

\begin{table}[htbp]
\centering
\small
\caption{Сравнение графовых моделей в warm-start сценарии. Таблица используется для проверки, какой способ учёта explicit feedback лучше работает для известных пользователей.}
\label{tab:graph_models}
\begin{tabular}{llrrr}
\toprule
Семейство & Модель & Warm HR@20 & Warm HR@50 & Warm nDCG@20 \\
\midrule
GCMC & \texttt{\detokenize{GCMC-rel}} & 0.1929 & 0.3123 & 0.0470 \\
GCMC & \texttt{\detokenize{GCMC}} & 0.1889 & 0.3121 & 0.0477 \\
GCMC & \texttt{\detokenize{GCMC-time}} & 0.1989 & 0.3280 & \textbf{0.0482} \\
NGCF & \texttt{\detokenize{NGCF-lin}} & 0.2024 & 0.3545 & 0.0452 \\
NGCF & \texttt{\detokenize{NGCF-exp}} & 0.2038 & 0.3555 & 0.0455 \\
NGCF & \texttt{\detokenize{NGCF-spline}} & \textbf{0.2108} & \textbf{0.3644} & 0.0473 \\
\bottomrule
\end{tabular}
\end{table}


Из таблицы~\ref{tab:graph_models} видно, что среди GCMC-вариантов наилучшее
качество показывает \texttt{GCMC-time}: $0.1989$ по HR@20 в режиме тёплого старта и $0.0482$ по
nDCG@20 в режиме тёплого старта. Это подтверждает вывод эксперимента
из подраздела~\ref{subsec:experiment_temporal_gcmc}: для GCMC временное
затухание является полезной модификацией.

Среди NGCF-вариантов лучший HR@20 показывает \texttt{NGCF-spline}:
$0.2108$ против $0.2024$ у \texttt{NGCF-lin} и $0.2038$ у \texttt{NGCF-exp}.
Это согласуется с результатами подраздела~\ref{subsec:experiment_rating_weight}:
обучаемая функция рейтингового веса оказывается полезнее фиксированных
монотонных функций. При этом по nDCG@20  \texttt{GCMC-time} немного
превосходит NGCF-варианты. Следовательно, разные способы учёта рейтинга
улучшают разные стороны ранжирования: обучаемый вес ребра лучше увеличивает
вероятность попадания релевантного объекта в первые $K$ позиций, а GCMC с типами рейтинговых связей
с учётом времени даёт немного более высокое качество верхних позиций списка.

Вывод эксперимента состоит в том, что графовые модели действительно позволяют
использовать явную обратную связь разными способами. Временное затухание улучшает
GCMC, а обучаемая функция рейтингового веса улучшает NGCF. Однако даже лучшие
графовые варианты в изолированном виде уступают двубашенной модели, что
проверяется далее при сравнении гибридных моделей и ансамблей.

\subsection{Гибридные модели и ансамбли}
\label{subsec:experiment_hybrid}

Основная проверяемая гипотеза работы состоит в том, что двубашенные и
графовые модели могут быть взаимодополняющими. Цель данного эксперимента состоит в том, чтобы
проверить, улучшается ли качество в режиме тёплого старта при передаче графовых
представлений в двубашенную модель или при ансамблировании оценок релевантности.
Двубашенная модель эффективно агрегирует историю пользователя и хорошо
подходит для сценария отбора кандидатов, а графовая модель явно использует структуру
графа пользователь-фильм и рейтинговые связи.

\begin{table}[htbp]
\centering
\small
\caption{Гибридные модели и ансамбли в warm-start сценарии. Таблица проверяет, даёт ли объединение двубашенной и графовой компонент выигрыш относительно базовой TT.}
\label{tab:hybrid_ensembles}
\begin{tabular}{lrrrr}
\toprule
Модель & Warm HR@20 & Warm HR@50 & Warm nDCG@20 & Warm nDCG@50 \\
\midrule
\texttt{\detokenize{TT}} & 0.3223 & 0.5128 & 0.0796 & 0.1140 \\
\texttt{\detokenize{TT<-GCMC}} & 0.3196 & 0.5128 & 0.0785 & 0.1121 \\
\texttt{\detokenize{TT<-NGCF}} & 0.1437 & 0.2629 & 0.0292 & 0.0442 \\
\texttt{\detokenize{Ens-TT-NGCF}} & 0.2208 & 0.3580 & 0.0542 & 0.0752 \\
\texttt{\detokenize{Ens-TT-GCMC-frz}} & 0.3251 & 0.5117 & 0.0806 & 0.1151 \\
\texttt{\detokenize{Ens-TT-GCMC-ft}} & \textbf{0.3432} & \textbf{0.5373} & \textbf{0.0874} & \textbf{0.1235} \\
\bottomrule
\end{tabular}
\end{table}


Из таблицы~\ref{tab:hybrid_ensembles} видно, что базовая двубашенная модель
\texttt{TT} существенно превосходит чистые графовые модели: она достигает
$0.3223$ по HR@20 и $0.0796$ по nDCG@20. Гибрид
\texttt{TT<-GCMC} даёт близкое качество: $0.3196$ по HR@20 и $0.0785$
по nDCG@20. Это означает, что передача GCMC-представлений в двубашенную
модель не разрушает качество, но сама по себе не даёт заметного выигрыша.

Наилучшее качество в режиме тёплого старта показывает ансамбль \texttt{Ens-TT-GCMC-ft}:
$0.3432$ по HR@20, $0.5373$ по HR@50 и $0.0874$ по nDCG@20.
Относительно базовой \texttt{TT} это даёт прирост примерно $6.5\%$ по
HR@20 и $9.8\%$ по nDCG@20. Следовательно, совместное использование
TT и GCMC действительно позволяет улучшить качество рекомендаций для известных
пользователей.

В таблице~\ref{tab:ensemble_alpha} приведены обученные коэффициенты
смешивания для суммы пользовательских представлений и представлений фильмов. В обоих вариантах
коэффициент $\alpha$ умножает двубашенную компоненту, а $1-\alpha$ умножает
графовую компоненту.

\begin{table}[H]
\centering
\caption{Коэффициенты смешивания в ансамбле TT и GCMC.}
\label{tab:ensemble_alpha}
\begin{tabular}{lcc}
\toprule
\textbf{Режим ансамбля} & $\boldsymbol{\alpha}_u$ & $\boldsymbol{\alpha}_i$ \\
\midrule
Замороженные базовые модели & $0.5598$ & $0.5327$ \\
Совместное дообучение ансамбля & $0.5171$ & $0.5161$ \\
\bottomrule
\end{tabular}
\end{table}

Значения коэффициентов близки к $0.5$, поэтому ансамбль не сводится к выбору
одной из двух моделей. В режиме с замороженными базовыми моделями
$\alpha_u=0.5598$ и $\alpha_i=0.5327$, то есть смешивание немного сильнее
опирается на двубашенную компоненту, особенно при построении пользовательского
представления. Это ожидаемо: TT напрямую агрегирует историю пользователя, тогда
как GCMC добавляет сигнал из локальной графовой структуры. После совместного
дообучения коэффициенты становятся ещё ближе к равному смешиванию
($\alpha_u=0.5171$, $\alpha_i=0.5161$). Это можно интерпретировать как
выравнивание пространств представлений: компоненты ансамбля подстраиваются друг
к другу, и ни одна из них не подавляет другую.

\begin{figure}[htbp]
    \centering
    \includegraphics[width=0.72\textwidth]{figures/validation_tt_vs_ens_gcmc.png}
    \caption{Сравнение качества \texttt{TT} и \texttt{Ens-TT-GCMC-ft}
    на валидационной выборке. Ансамбль достигает более высокого HR@20,
    что подтверждает взаимодополняемость двубашенной и графовой компонент.}
    \label{fig:валидации_tt_vs_ens_gcmc}
\end{figure}

График~\ref{fig:валидации_tt_vs_ens_gcmc} показывает, что преимущество
ансамбля наблюдается на валидационной выборке уже с первых эпох.

Вывод эксперимента состоит в том, что ансамблирование двубашенной и GCMC-модели
является наиболее полезной модификацией для сценария тёплого старта. Объединение предобученных моделей на небольшом количестве эпох, при объединении будет давать прирост качества. При этом
устойчивость таких ансамблей к новым пользователям требует отдельной проверки,
которая проводится в следующем подразделе.

\subsection{Отдельная проверка сценария холодного старта по пользователю}
\label{subsec:experiment_cold_user}

Отдельный эксперимент посвящён холодному старту по пользователям. Его цель состоит в том, чтобы
проверить, какие модели сохраняют качество, когда пользователь отсутствует в
основной обучающей выборке и его представление нужно построить только по
опорной истории. Это важно для практического применения, поскольку в
рекомендательной системе постоянно появляются новые пользователи.

\begin{table}[htbp]
\centering
\small
\caption{Сравнение моделей в cold-user сценарии. Для новых пользователей качество измеряется после построения пользовательского представления по support-истории.}
\label{tab:cold_user_results}
\begin{tabular}{llrrrr}
\toprule
Семейство & Модель & HR@20 & HR@50 & nDCG@20 & nDCG@50 \\
\midrule
Two-Tower & \texttt{\detokenize{TT}} & 0.3156 & 0.5057 & 0.0778 & 0.1109 \\
Two-Tower & \texttt{\detokenize{TT-bin}} & \textbf{0.3239} & \textbf{0.5071} & \textbf{0.0806} & \textbf{0.1123} \\
GCMC & \texttt{\detokenize{GCMC}} & 0.1892 & 0.3062 & 0.0470 & 0.0641 \\
GCMC & \texttt{\detokenize{GCMC-time}} & 0.1946 & 0.3134 & 0.0469 & 0.0645 \\
NGCF & \texttt{\detokenize{NGCF-lin}} & 0.0994 & 0.1649 & 0.0193 & 0.0268 \\
NGCF & \texttt{\detokenize{NGCF-exp}} & 0.0999 & 0.1651 & 0.0194 & 0.0268 \\
NGCF & \texttt{\detokenize{NGCF-spline}} & 0.0003 & 0.0004 & 0.0000 & 0.0001 \\
Hybrid & \texttt{\detokenize{TT<-GCMC}} & 0.3171 & 0.5034 & 0.0779 & 0.1099 \\
Hybrid & \texttt{\detokenize{TT<-NGCF}} & 0.1431 & 0.2558 & 0.0293 & 0.0436 \\
Ensemble & \texttt{\detokenize{Ens-TT-GCMC-frz}} & 0.2825 & 0.4664 & 0.0667 & 0.0982 \\
Ensemble & \texttt{\detokenize{Ens-TT-GCMC-ft}} & 0.2603 & 0.4280 & 0.0618 & 0.0901 \\
\bottomrule
\end{tabular}
\end{table}


Таблица~\ref{tab:cold_user_results} показывает, что лучшей моделью в
сценарии холодного старта по пользователю является \texttt{TT-bin}: она достигает $0.3239$ по HR@20,
$0.5071$ по HR@50 и $0.0806$ по nDCG@20. Базовая \texttt{TT} также сохраняет
высокое качество: $0.3156$ по HR@20 и $0.0778$ по nDCG@20. Гибрид
\texttt{TT<-GCMC} находится между ними и даёт $0.3171$ по HR@20.
Следовательно, архитектуры, которые строят пользовательское представление по
истории, оказываются наиболее устойчивыми для новых пользователей.

Графовые модели в данном сценарии ведут себя по-разному. \texttt{GCMC-time}
немного превосходит базовую \texttt{GCMC} по HR@20 ($0.1946$ против $0.1892$),
что согласуется с экспериментом по временному затуханию. Однако NGCF-варианты
существенно уступают TT-моделям, а \texttt{NGCF-spline} практически теряет
качество. Это показывает, что обучаемая функция рейтингового веса хорошо
работает для известных пользователей, но в текущей постановке плохо переносится
на новые опорные истории.

Ансамбль \texttt{Ens-TT-GCMC-ft}, который был лучшим в сценарии тёплого старта,
теперь снижается до $0.2603$ по HR@20. Это хуже, чем у базовой
\texttt{TT} и \texttt{TT-bin}. Возможное объяснение состоит в том, что
совместное дообучение ансамбля усиливает подстройку под распределение тёплого старта и
ухудшает перенос на новых пользователей. Поэтому максимизация качества в режиме тёплого старта и устойчивость к сценарию холодного старта по пользователю являются разными целями.

\begin{figure}[h]
    \centering
    \includegraphics[width=0.72\textwidth]{figures/warm_vs_cold_user_tradeoff}
    \caption{Сравнение HR@20 в режиме тёплого старта и HR@20 для холодного старта по пользователю. Модель
    \texttt{TT-bin} оказывается наиболее устойчивой к сценарию холодного старта по пользователю,
    тогда как \texttt{Ens-TT-GCMC-ft} максимизирует качество на известных
    пользователях.}
    \label{fig:warm_vs_cold_user_tradeoff}
\end{figure}

Вывод эксперимента состоит в том, что для новых пользователей наиболее
надёжной оказывается двубашенная модель, с учётом рейтинга только в качестве положительного сигнала. Ансамбль TT и GCMC полезен
для известных пользователей, но без дополнительной регуляризации или отдельной
настройки он хуже переносится на новых пользователей.

\subsection{Итоговое сравнение качества и времени}
\label{subsec:experiment_final}

Итоговый эксперимент объединяет основные результаты и сравнивает не только
качество, но и время оценки. Его цель состоит в том, чтобы показать, какие модели являются
предпочтительными в разных практических режимах: максимизация качества в режиме тёплого старта, устойчивость к новым пользователям или более быстрая оценка.
В таблицу~\ref{tab:final_quality_time} включены ключевые представители
семейств: матричная факторизация, двубашенная модель, бинарная двубашенная
модель, лучшие графовые варианты, гибридная модель и лучший ансамбль.


\begin{figure}[h]
    \centering
    \includegraphics[width=0.72\textwidth]{figures/quality_time_tradeoff}
    \caption{Компромисс между качеством HR@20 в режиме тёплого старта и временем оценки на
    $100\,000$ пользователей. Ансамбль \texttt{Ens-TT-GCMC-ft} показывает
    лучшее качество в режиме тёплого старта, но простые графовые модели быстрее.}
    \label{fig:quality_time_tradeoff}
\end{figure}

\begin{table}[htbp]
\centering
\small
\caption{Итоговое сравнение качества и времени оценки. В таблицу включены ключевые представители семейств моделей.}
\label{tab:final_quality_time}
\begin{tabular}{lrrrrr}
\toprule
Модель & Warm HR@20 & Warm nDCG@20 & Cold HR@20 & Cold nDCG@20 & Время на 100k, сек \\
\midrule
\texttt{\detokenize{MF-128}} & 0.2458 & 0.0616 & --- & --- & 71.52 \\
\texttt{\detokenize{TT}} & 0.3223 & 0.0796 & 0.3156 & 0.0778 & 90.48 \\
\texttt{\detokenize{TT-bin}} & 0.3207 & 0.0792 & \textbf{0.3239} & \textbf{0.0806} & 89.70 \\
\texttt{\detokenize{GCMC-time}} & 0.1989 & 0.0482 & 0.1946 & 0.0469 & 69.10 \\
\texttt{\detokenize{NGCF-spline}} & 0.2108 & 0.0473 & 0.0003 & 0.0000 & \textbf{65.67} \\
\texttt{\detokenize{TT<-GCMC}} & 0.3196 & 0.0785 & 0.3171 & 0.0779 & 94.98 \\
\texttt{\detokenize{Ens-TT-GCMC-ft}} & \textbf{0.3432} & \textbf{0.0874} & 0.2603 & 0.0618 & 83.27 \\
\bottomrule
\end{tabular}
\end{table}


Таблица~\ref{tab:final_quality_time} и рис.~\ref{fig:quality_time_tradeoff}
показывают, что разные модели оказываются оптимальными в разных режимах.
Если основной целью является максимизация качества на известных пользователях,
лучшим вариантом является \texttt{Ens-TT-GCMC-ft}: он достигает $0.3432$ по
HR@20 в режиме тёплого старта и $0.0874$ по nDCG@20 в режиме тёплого старта. Относительно сильной одиночной
двубашенной модели \texttt{TT} это соответствует приросту с $0.3223$ до
$0.3432$ по HR@20 в режиме тёплого старта, то есть примерно на $6.5\%$, и с $0.0796$ до
$0.0874$ по nDCG@20 в режиме тёплого старта, то есть примерно на $9.8\%$. Качественно этот
эффект можно объяснить тем, что TT и GCMC используют разные источники сигнала:
TT лучше агрегирует историю пользователя и признаки объектов, а GCMC явно
использует локальную структуру графа взаимодействий и рейтинговые связи.
Ошибки этих моделей оказываются не полностью совпадающими, поэтому их сумма
даёт более устойчивую оценку релевантности для ранжирования в режиме тёплого старта.

Если же важна устойчивость к новым пользователям, лучшей оказывается
\texttt{TT-bin}: она достигает $0.3239$ по HR@20  и $0.0806$ по
nDCG@20. Это означает, что выигрыш ансамбля в режиме тёплого старта не
переносится автоматически на постановку холодного старта по пользователю: для новых пользователей
важнее способ построения представления по опорной истории, чем дополнительный
графовая оценка релевантности.

С точки зрения времени оценки наиболее быстрым вариантом среди моделей
итогового сравнения является \texttt{NGCF-spline}: $65.67$ секунды на
$100\,000$ пользователей. Однако его качество в режиме холодного старта по пользователю практически нулевое,
поэтому быстрота оценки не компенсирует потерю устойчивости. Модель \texttt{GCMC-time}
также работает быстрее двубашенных моделей, но уступает им по качеству.
Следовательно, для практического применения требуется учитывать не только
задержку ответа, но и качество в разных сценариях.


\section{Заключение}
\label{sec:conclusion}

В работе исследовалось применение графовых нейронных сетей в задаче построения списка из первых $K$
рекомендаций фильмов по истории явных пользовательских оценок. Основной акцент
был сделан на использовании графа пользователь-фильм и значений явных
рейтингов для построения более информативных представлений пользователей и
объектов.

В ходе работы был подготовлен датасет с дополнительными признаками фильмов,
временным разбиением и отдельными сценариями тёплого старта и холодного старта по пользователю. Сохранение
исходных численных рейтингов позволило рассматривать несколько способов их
использования: бинарную релевантность, рейтинговые веса, типы рёбер и
временно-взвешенные связи. В экспериментах сравнивались двубашенная модель как
сильное неграфовое базовое решение, GCMC, NGCF-подобные графовые модели, а также
гибридные модели и ансамбли двубашенной модели с GCMC.

Эксперименты показали, что способ задания рёбер в пользователь--фильм графе заметно
влияет на качество рекомендаций. Для GCMC полезным оказалось временное
затухание, так как более свежие оценки лучше отражают текущие предпочтения
пользователя. Для NGCF-подобной модели лучший результат среди вариантов с
рейтинговыми весами дала обучаемая функция веса, что подтверждает важность
выбора зависимости между значением рейтинга и силой связи в графе.

Лучший результат в сценарии тёплого старта был получен ансамблем двубашенной
SASRec-подобной модели и GCMC. По сравнению с одиночной двубашенной моделью
ансамбль улучшил HR@20 на $6.5\%$, а nDCG@20 на $9.8\%$. Это показывает,
что графовая компонента добавляет структурный сигнал о связях пользователей и
фильмов, который не полностью извлекается неграфовой моделью.

Коэффициенты смешивания также подтверждают дополняющий характер моделей. При
замороженных базовых моделях получены $\alpha_i=0.5327$ и $\alpha_u=0.5598$,
а при совместном дообучении получены $\alpha_i=0.5161$ и $\alpha_u=0.5171$. Эти
значения близки к $0.5$, поэтому ансамбль не сводится к доминированию одной
архитектуры, а использует почти равный вклад двубашенного и графового
представлений.

Основным ограничением остаётся различие между сценариями тёплого старта и холодного старта по пользователю. Ансамбль TT и GCMC улучшает рекомендации для известных
пользователей, но хуже переносится на новых пользователей, чем бинарная
двубашенная модель \texttt{TT-bin}. Это связано с тем, что графовые модели
наиболее эффективны для пользователей, уже присутствующих в обучающем графе,
и требуют дополнительных механизмов для построения представлений новых
пользователей.

Таким образом, полученные результаты показывают перспективность дальнейшей
разработки гибридных рекомендательных моделей, объединяющих графовые и двубашенные нейросети, но с возможной регуляризацией для новых пользователей. Кроме того, эксперименты подтверждают, что при наличии явных оценок их
целесообразно учитывать не только как факт взаимодействия, но и как
дополнительный сигнал о силе пользовательского предпочтения.

\clearpage
\renewcommand{\refname}{Список литературы}
\addcontentsline{toc}{section}{Список литературы}
\bibliographystyle{ugost2008}
\bibliography{references_recommender}

\end{document}
